\documentclass[11pt]{article}
\usepackage{amsmath,amssymb,amsthm,hyperref,geometry}\geometry{margin=1in}
\newtheorem{theorem}{Theorem}\newtheorem{proposition}{Proposition}
\title{Multicategorical coupling wiring and dependent context interchange}
\date{15 November 2026}\begin{document}\maketitle
\paragraph{Status.} Scoped abstract construction with written proofs. The general construction uses a category, then specializes to the actual effect-fibre contextad. It retains chosen dependence witnesses as operations, rather than asserting a natural or associative selection from marginal effects. It supplies hierarchical wiring and a compatibility law with dependent context composition. It is not a claim that arbitrary overlapping coupling requirements admit a joint extension or that a complete concurrent interactad has been constructed.

\section{A symmetric multicategory of explicit cones}
This elementary opposite-category cone construction is standard categorical bookkeeping, not a novelty claim. Its role here is to expose which joint data are retained and to state the subsequent native dependent-context interchange precisely.
For any locally small category $\mathcal E$, define a symmetric multicategory $\mathsf{Cone}(\mathcal E)$ with the same objects and
\[
 \mathsf{Cone}(\mathcal E)(P_1,\ldots,P_n;Q)
   =\prod_{i=1}^n\mathcal E(Q,P_i).
\]
The output $Q$ is the retained joint object; the input objects are its projected components. The unary category is $\mathcal E^{\mathrm{op}}$. Substitution of $u_i:Q\to P_i$ and $v_{ij}:P_i\to P_{ij}$ is the family $v_{ij}u_i:Q\to P_{ij}$. Identities and symmetric actions are identities and permutations of projection lists. The nullary operation at $Q$ is the empty family, not a chosen map from a terminal effect object.

\begin{proposition}[Coherent hierarchical wiring]
These operations satisfy the symmetric multicategory laws. A functor $F:\mathcal E\to\mathcal E'$ induces a multifunctor by applying $F$ to every retained object and projection. In particular a tree of specified local cone operations has a composite whose root object is unchanged and whose leaf projections are the corresponding composites; every order of performing substitutions gives the same native projections.
\end{proposition}
\begin{proof}
Both substitution orders send a leaf along the same composable chain of arrows of $\mathcal E$. Associativity and units in $\mathcal E$ give equality. Permutations only reorder labels and commute with those composites. Functoriality follows from preservation of those operations by $F$. Nothing in this proof identifies two different root objects or two different joint effects with the same projections.
\end{proof}

For $\mathcal E_T$ the category of supplied effects $(X,p\in TX)$ and deterministic effect-preserving maps, an operation is exactly an actual joint effect presentation with native marginal projections. The given root joint is local to one wiring node; recursive substitution derives its projections to distant leaves. One is not asked to provide the final leaf-level coupling equations separately. All chosen intermediate joint presentations and their maps may be retained in a tree presentation; evaluating the tree composes maps but does not erase its root witness.

\section{Why this is not a hidden canonical tensor}
\begin{proposition}[Representability boundary]
This multicategory is representable by a monoidal structure on its underlying unary category precisely when $\mathcal E$ has finite products representing the cone functors. In that case its tensor objects are those products, viewed in $\mathcal E^{\mathrm{op}}$.
\end{proposition}
\begin{proof}
A representing tensor for $(P_i)$ is by definition an object $R$ with natural bijections
in the unary orientation $\mathcal E^{\mathrm{op}}(R,Q)=\mathcal E(Q,R)$, explicitly
\[
 \mathcal E(Q,R)\cong\prod_i\mathcal E(Q,P_i)
\]
for every $Q$. This is exactly the universal property of a product. The nullary case is a terminal object. Conversely finite products give these bijections and their standard coherent monoidal comparisons.
\end{proof}

For finite normalized distributions with deterministic measure-preserving maps, two fair-bit objects do not have such a product. If a universal object $R$ with projections existed, pushing its effect along the projection pair would give one fixed joint law $\sigma$. The independent fair-bit coupling and the diagonal fair-bit coupling are both cones. Their unique mediating maps to $R$ would be effect-preserving and commute with projections, forcing $\sigma$ to equal each of their different joint laws, a contradiction. Thus no representing tensor can cover all such dependence witnesses. This does not rule out a chosen independent tensor, which lacks that universal cone property.

Likewise for nonempty-powerset effects with image-preserving maps, the full allowed relation on two bits and the diagonal allowed relation would force incompatible allowed sets as the product object's projected image. These are two nonprobabilistic dependence witnesses in the same abstract wiring construction.

\section{Actual effect-context instantiation and pure reindexing}
For every pure set $A$, the effect-fibre context category $\mathcal K_A$ has objects $(X_a,p_a)$ and pointwise effect-preserving maps. Pure reindexing $f^*:\mathcal K_{A'}\to\mathcal K_A$ therefore induces a multifunctor $\mathsf{Cone}(\mathcal K_{A'})\to\mathsf{Cone}(\mathcal K_A)$. Its identities and composite-map laws are the original context reindexing laws, on actual native maps. If a class $\mathcal J$ of admitted carrier projections is used, it must be closed as in the context construction; cone maps themselves need not lie in $\mathcal J$.

The next result relates this wiring to dependent composition rather than only giving an unrelated category of marginals. Let a first cone have projections $u_i:Z\to P_i$ over $A$. Let $R$ be a context over $e_AZ$, and let $R_i$ be contexts over $e_AP_i$. Supply native conditional maps
\[
 v_i:R\longrightarrow(e_Au_i)^*R_i.
\tag{L}
\]
The object $R$ and its effect at each actual joint witness may depend on that entire witness. Equation (L) requires its projected effect to be the declared continuation at the projected component witness.

\begin{theorem}[Dependent serial interchange with cone substitution]
Dependent context composition gives a cone
\[
 Z\star R\longrightarrow P_i\star R_i,
 \qquad(z,r)\longmapsto(u_i z,v_{i,z}r).
\tag{A}
\]
It is native and compatible with pure reindexing, native cone maps, context units/associators and further hierarchical cone substitution. In particular composing local projections before or after dependent serial assembly gives identical raw maps, effects and context extension maps.
\end{theorem}
\begin{proof}
The effect equation is bind naturality: after (L), the inner projected effect depends on $z$ only through $u_i z$, allowing the outer effect to be pushed along $u_i$. The dependent-sum carrier is precisely the original context multiplication, so no marginal-only equivalence is used.

For the nontrivial substitution check, suppose further projections $t_{ij}:P_i\to P_{ij}$ and continuation maps $s_{ij}:R_i\to(e_At_{ij})^*R_{ij}$ are supplied. The two routes both send
\[
 (z,r)\longmapsto
 \bigl(t_{ij}u_i z,\ s_{ij,u_i z}(v_{i,z}r)\bigr).
\]
They therefore agree on complete raw data; their effect equations are the composite native equations. Here naturality in cone data means compatible maps of both levels: $a:Z\to Z\prime$ and $b:R\to(e_Aa)^*R\prime$, with $u_i=u_i\prime a$ and $v_{i,z}=v_{i,a(z)}\prime b_z$, induce $(z,r)\mapsto(az,b_zr)$. A map of the first context alone is insufficient. The same coordinate calculation with triples proves compatibility with the context associator. Unit comparisons remove only unit singleton coordinates, using monad unit laws. Reindexing replaces $a$ by $f(a)$ everywhere. All further diagrams are compositions of these equalities, not isomorphisms between arbitrarily selected couplings.
\end{proof}

Adding pure bodies and the matching interface-projection squares yields the operation-level conditional assembly theorem in the companion note. Its essential factorization is exactly (L). When a crossing already acts functorially on the relevant native context/operation categories, the first proposition transports these explicit cones through that functor. This observation preserves existing native naturality; it does not independently prove a new crossing for a class in which one has not been constructed.

\section{Local existence versus global substitution}
The companion admission criterion supplies local cone objects when the effect functor weakly preserves the appropriate pullback. Finite-distribution shared-boundary equality and nonempty-powerset matching-image equality instantiate that property with explicit witnesses. A tree of those admitted operations can then be substituted by the theorem above. No coherent section of the weak-pullback comparison map is assumed: the selected joint objects are retained as node data, and substituting the same data is coherent.

Different choices of node data can produce different dependence. A cyclic overlap problem is not a tree of freely substitutable native projections unless a suitable joint witness has already been admitted; pairwise marginal consistency alone is not asserted to solve it. This is the precise difference between the positive hierarchical assembly theorem and an unsupported universal joint-extension claim.
\end{document}
